\section{Yang--Mills as a full-chain test}
\label{sec:ymchain}

The Yang--Mills problem is unusually revealing because the structures needed
for the conclusion are naturally distributed across representations.  The
official target combines existence of a nontrivial continuum quantum field
theory with a positive Hamiltonian mass gap \cite{JaffeWitten2000}.  A finite
lattice construction makes neither conjunct automatic.

\subsection{The representation chain}

Fix a compact simple gauge group $G$.  A constructive route may begin with
lattice link variables at spacing $a$ and physical volume $L$, pass to
gauge-orbit data, integrate fluctuations into block variables, take volume
and cutoff limits of renormalized gauge-invariant observables, reconstruct a
Hilbert-space theory, and finally interpret Euclidean decay spectrally:

\begin{figure}[t]
\centering
\begin{tikzpicture}[
  node distance=7mm,
  box/.style={draw,rounded corners,align=center,minimum width=42mm,
    minimum height=8mm,font=\small},
  arr/.style={-{Latex[length=1.8mm]},thick},
  lab/.style={font=\scriptsize,align=left},
  exposed/.style={font=\scriptsize,align=left,text width=47mm}]
\node[box] (lat) {lattice links};
\node[box,below=of lat] (quo) {gauge quotient};
\node[box,below=of quo] (rg) {blocked family};
\node[box,below=of rg] (sch) {continuum Schwinger data};
\node[box,below=of sch] (os) {OS Hilbert space};
\node[box,below=of os] (sp) {Hamiltonian spectrum};
\draw[arr] (lat) -- node[left,lab] {descent} (quo);
\draw[arr] (quo) -- node[left,lab] {block RG} (rg);
\draw[arr] (rg) -- node[left,lab] {limits} (sch);
\draw[arr] (sch) -- node[left,lab] {OS reconstruction} (os);
\draw[arr] (os) -- node[left,lab] {spectral theorem} (sp);
\node[exposed,right=15mm of lat] {manifest structure: locality};
\node[exposed,right=15mm of quo] {manifest structure: gauge invariance};
\node[exposed,right=15mm of rg] {manifest structure: scale dependence};
\node[exposed,right=15mm of sch] {manifest structure: reflection positivity};
\node[exposed,right=15mm of os] {manifest structure: physical time};
\node[exposed,right=15mm of sp] {manifest structure: spectral gap};
\end{tikzpicture}
\caption{The Yang--Mills representation chain.  The labels to the right of the nodes
name structures naturally exposed there, not properties absent everywhere
else.  Every arrow has separate preservation and recovery obligations.}
\label{fig:ymchain}
\end{figure}

The chain is not a chronology and no node is globally superior.  It is an
interface graph.  The main obligations are as follows.

\paragraph{Gauge descent.}
The quotient must preserve the physical observable algebra and measure while
accounting for orbit singularities, residual gauge transformations, and
topological sectors.  A local gauge chart can be exact on its domain without
covering the full measure.  Conversely, unreduced variables can carry useful
local structure even when they are redundant; this is consonant with the
locality analysis of surplus gauge structure
\cite{NguyenTehWells2020}.

\paragraph{Scale transformation.}
An RG step changes variables and often the form or range of the action.  Its
physical meaning depends on the specific RG construction; coarse graining is
not synonymous with every use of the renormalization group
\cite{KoberinskiFraser2023}.  For a constructive proof, the map must retain a
state sufficient to reconstruct the required observables and must control
positivity, locality range, large fields, boundary conditions, and topology.

\paragraph{Infinite volume and continuum.}
One needs normalized gauge-invariant observables whose correlation
functionals converge in a declared distributional topology, together with
reflection positivity, Euclidean covariance, clustering or another vacuum
selection condition, nontriviality, and regulator independence.  Constants
must be uniform in $L$ before $L\to\infty$ and stable in physical units as
$a\downarrow0$.  A subsequential symmetric limit is not automatically the
intended pure Yang--Mills theory.

\paragraph{Reconstruction and spectral interpretation.}
The OS axioms provide a theorem-backed route from suitable Euclidean data to
a Hilbert space and positive Hamiltonian
\cite{OsterwalderSchrader1973,OsterwalderSchrader1975}.  To infer a global gap
from decay, the centered gauge-invariant core used in the estimate must be
dense in the reconstructed vacuum complement.  Theorem
\ref{thm:recovery-spectrum} then performs only the final inference; it cannot
supply the Euclidean estimate or the continuum object.

\subsection{Property visibility}

The decisive properties do not line up in one column.  Lattice variables make
local action and exact finite gauge covariance explicit.  The quotient makes
physical invariance explicit but can hide local sections.  RG variables make
scale dependence explicit but can hide microscopic locality and normalized
boundary data.  Schwinger functions make Euclidean correlation structure
explicit, while a Hamiltonian makes spectrum explicit.  Reflection
positivity is an action- and interface-specific certificate, not an automatic
benefit of Euclidean notation.  Appendix~\ref{app:ymtable} supplies the full
visibility and obligation matrix.

\subsection{The exact first open physical implication}

The chain isolates a frontier before OS reconstruction.  One needs a
gauge-invariant observable core $\mathcal D_{a,L}$ and constants in physical
units such that, schematically,
\begin{equation}
\label{eq:ym-rigidity}
 \left|\langle A\,\tau_t(B)\rangle_{a,L,c}\right|
 \leq C_{A,B}\e^{-mt}+r_{a,L}(A,B,t),
\end{equation}
with $m>0$ and admissible prefactors uniform in volume and sufficiently small
cutoff, and with recovery errors converging in the topology used for the
Schwinger limit.  The statement must cover a core that becomes dense after
reconstruction.  It may be generated by RG contraction, a polymer expansion,
a transfer estimate, or another noncircular rigidity engine.  Assuming
all-distance exponential clustering as an input would merely assume the
needed property in a different representation.

No result in this paper proves Equation~\eqref{eq:ym-rigidity} for
four-dimensional weak-coupling pure Yang--Mills, nor does it prove the
continuum premises.  The contribution is the implication boundary: local or
finite spectral rigidity counts as a physical mass gap only after the
full-measure, uniform, continuum, positivity, completeness, and reconstruction
certificates close.
